% Standalone problem document, generated from the adjacent README.md.
% Compile with XeLaTeX. Mathematical content is maintained in Markdown.
\documentclass[11pt,a4paper]{article}
\usepackage[fontsize=10.5pt]{scrextend}
\usepackage[margin=22mm,headheight=15pt,headsep=9mm,footskip=11mm]{geometry}
\usepackage{fontspec}
\setmainfont{texgyrepagella}[Extension=.otf,UprightFont=*-regular,BoldFont=*-bold,ItalicFont=*-italic,BoldItalicFont=*-bolditalic]
\setsansfont{texgyreheros}[Extension=.otf,UprightFont=*-regular,BoldFont=*-bold,ItalicFont=*-italic,BoldItalicFont=*-bolditalic]
\setmonofont{texgyrecursor}[Extension=.otf,UprightFont=*-regular,BoldFont=*-bold,ItalicFont=*-italic,BoldItalicFont=*-bolditalic,Scale=MatchLowercase]
\usepackage{amsmath,amssymb}
\usepackage{unicode-math}
\setmathfont{texgyrepagella-math.otf}
\usepackage{xcolor,microtype,enumitem,needspace,fancyhdr}
\usepackage{bookmark}
\definecolor{ink}{HTML}{17344B}
\definecolor{accent}{HTML}{197D80}
\definecolor{muted}{HTML}{596773}
\hypersetup{colorlinks=true,linkcolor=accent,urlcolor=accent,pdftitle={IE-02 - Is
the ideal GMRES bound sharp for every Jordan
block?},pdfauthor={Open Problems in Numerical Linear Algebra}}
\urlstyle{same}
\setlength{\parindent}{0pt}
\setlength{\parskip}{5pt plus 1pt minus 1pt}
\linespread{1.04}
\setlength{\emergencystretch}{3em}
\widowpenalty=10000
\clubpenalty=10000
\displaywidowpenalty=10000
\allowdisplaybreaks[1]
\setlist{nosep,leftmargin=1.5em}
\providecommand{\tightlist}{\setlength{\itemsep}{0pt}\setlength{\parskip}{0pt}}
\providecommand{\pandocbounded}[1]{#1}
\pagestyle{fancy}
\fancyhf{}
\fancyhead[L]{\sffamily\footnotesize\color{muted}OPEN PROBLEMS IN NUMERICAL LINEAR ALGEBRA}
\fancyhead[R]{\sffamily\bfseries\footnotesize\color{accent}IE-02}
\fancyfoot[L]{\parbox[b]{0.9\textwidth}{\sffamily\fontsize{8}{10}\selectfont\color{muted}Editorial ratings; a bounded search cannot certify that a problem remains unsolved.\\Literature check: 2026-09-11}}
\fancyfoot[R]{\sffamily\footnotesize\color{muted}\thepage}
\renewcommand{\headrulewidth}{0.3pt}
\renewcommand{\headrule}{\hbox to\headwidth{\color{accent}\leaders\hrule height \headrulewidth\hfill}}
\makeatletter
\renewcommand{\section}{\Needspace{5\baselineskip}\@startsection{section}{1}{0pt}{12pt plus 2pt minus 2pt}{4pt}{\sffamily\bfseries\large\color{ink}}}
\renewcommand{\subsection}{\Needspace{4\baselineskip}\@startsection{subsection}{2}{0pt}{9pt plus 1pt minus 1pt}{3pt}{\sffamily\bfseries\normalsize\color{ink}}}
\makeatother
\setcounter{secnumdepth}{0}
\begin{document}
{\sffamily\small\color{muted}Linear systems and elimination\par}
\vspace{3pt}
{\raggedright\hyphenpenalty=10000\exhyphenpenalty=10000\sffamily\bfseries\fontsize{21}{24}\selectfont\color{ink}Is
the ideal GMRES bound sharp for every Jordan block?\par}
\vspace{7pt}
{\sffamily\small\textbf{Difficulty:} challenging\quad\textbf{Importance:} interesting
to specialist\par}
{\sffamily\small\bfseries\color{accent}Status: Solved\par}
\vspace{4pt}
{\color{accent}\hrule height 0.8pt}
\vspace{6pt}
\textbf{Rating rationale:} Historical assessment retained: challenging
because the remaining Jordan-block minimax cases require control of
multiple extremal singular vectors; specialist impact reflects a
structural test case for GMRES theory.

\subsection{Resolution --- 2026-09-11}\label{resolution-2026-09-11}

\textbf{Solved affirmatively.} George Stepaniants, Department of
Computing and Mathematical Sciences, California Institute of Technology,
Pasadena, California, USA, proves the displayed equality for every
\(n\ge2\), \(1\le k<n\), and nonzero complex \(\lambda\), with complex
polynomials and starting vectors. No divisibility or eigenvalue-regime
case remains open.

\href{https://github.com/ajt60gaibb/OpenProblemsInNLA/blob/main/linear-systems-and-elimination/IE-02/solution.md}{Theorem
1 and its proof} establish the exact original target. The stronger
Theorem 6 proves affine minimax equality for triangular Toeplitz
matrices: finite Carathéodory--Fejér interpolation describes the maximal
singular subspace, and scalar spectral factorization preserves every
complex GMRES orthogonality equation in one unit vector.
\href{https://github.com/ajt60gaibb/OpenProblemsInNLA/blob/main/linear-systems-and-elimination/IE-02/solution.pdf}{Proof
PDF} ·
\href{https://github.com/ajt60gaibb/OpenProblemsInNLA/blob/main/linear-systems-and-elimination/IE-02/solution.tex}{Standalone
XeLaTeX source}.

The complete AI-assisted proof passed a separate
\href{https://github.com/ajt60gaibb/OpenProblemsInNLA/blob/main/references/stepaniants-ie02-2026-09-11/verification/IE-02-independent-review.md}{independent
Codex-agent review}. This is automated-agent verification, without a
claim of external human peer review or formal certification.
\href{https://github.com/ajt60gaibb/OpenProblemsInNLA/blob/main/references/stepaniants-ie02-2026-09-11/README.md}{Submission
record, preserved source and public eligibility check}. The original
problem statement and permanent ID remain unchanged; the earlier special
cases and status checks below are retained as history.

\subsection{Problem statement}\label{problem-statement}

For \(n\geq2\) and \(\lambda\in\mathbb C\setminus\{0\}\), let
\(J_n(\lambda)=\lambda I+N\), where \(N_{i,i+1}=1\) and all other
entries of \(N\) vanish. Let
\(\mathcal P_k=\{p\in\mathbb C[z]:\deg p\leq k,\ p(0)=1\}\). Define

\[
\psi_k(J)=\max_{\|v\|_2=1}\min_{p\in\mathcal P_k}\|p(J)v\|_2,
\qquad
\phi_k(J)=\min_{p\in\mathcal P_k}\|p(J)\|_2.
\]

Prove or disprove \(\psi_k(J_n(\lambda))=\phi_k(J_n(\lambda))\) for
every \(1\leq k<n\). The maximum describes the slowest possible GMRES
residual reduction, whereas the minimum over operator norms is the ideal
bound. The familiar inequality \(\psi_k\leq\phi_k\) does not answer the
question. A single Jordan block is a published structural test case, not
a claim about all nonnormal matrices.

\newpage
\subsection{References}\label{references}

Tichý, Liesen, and Faber,
\href{https://etna.ricam.oeaw.ac.at/volumes/2001-2010/vol26/abstract.php?pages=453-473}{\emph{On
worst-case GMRES, ideal GMRES, and the polynomial numerical hull of a
Jordan block}}, ETNA 26 (2007), 453--473, §1 conjecture and subsequent
special cases. Faber, Liesen, and Tichý,
\href{https://arxiv.org/abs/2506.09687}{\emph{Matrix best approximation
in the spectral norm}}, published in LAA 733 (2026), §§4--5.

\subsection{Earlier status check ---
2026-09-08}\label{earlier-status-check-2026-09-08}

Searches for
\textit{Jordan block ideal GMRES equality proved 2026} and
\textit{site:arxiv.org GMRES Jordan block} found the original partial
results and the 2026 general approximation paper, but no resolution of
the displayed Jordan-block equality. The latter paper's doubling theorem
changes the matrix and does not by itself establish this statement.

\subsection{Audit update --- 2026-09-10}\label{audit-update-2026-09-10}

Rechecked the
\href{https://www.karlin.mff.cuni.cz/~ptichy/download/public/TiLiFa2007.pdf}{author
copy of Tichý--Liesen--Faber}, especially §§3--5: Corollary 4.4 proves
equality whenever \(k\) divides \(n\), with additional eigenvalue
regimes proved elsewhere in those sections. These are substantive parts
of the displayed target. The
\href{https://arxiv.org/html/2506.09687}{2025/2026 approximation paper}
and targeted Jordan-block/ideal-GMRES searches did not supply the
remaining cases; its doubling construction changes the input matrix.
\end{document}
